\subsection{Formalization of Motion Primitives}\label{sec:motion-primitives}

As motivated in \Cref{sec:introduction}, there exist a variety of prehensile and nonprehensile motion primitives, each with its core set of capabilities and limitations.
We unify this diverse portfolio under a single definition of \textsl{motion primitive} $MP$ as an entity with four components: 
i) $f_{MP}$, a forward model (\Cref{def:forward-model}); 
ii) $\phi_{MP}$, a control mechanism for generating the robot--object motion (\Cref{def:control-mechanism}); 
iii) $\mathcal{M}_{MP}$, a situated kinodynamic manifold (\Cref{def:situated-manifold}); and, 
iv) $proj_{MP}$, a control projection function (\Cref{def:projection}).
A more extensive treatment of these components and our design choices is presented in the context of our motion planning framework in \Cref{sec:methods-planning}. 

% \begin{enumerate}
%     \item $f_{MP}$, a forward model that computes the final object pose from an initial pose and a set of control parameters (\Cref{def:forward-model});
%     \item $\phi_{MP}$, a control mechanism for generating the robot--object motion resulting from a set of control parameters (\Cref{def:control-mechanism});
%     \item $\mathcal{M}_{MP}$, a situated kinodynamic manifold that maps the current environment state and primitive control parameters to a real value to determine goodness of action (\Cref{def:situated-manifold}); and,
%     \item $proj_{MP}$, a control projection function that maps a generic robot dynamics rollout to the primitive's situated kinodynamic manifold, thus allowing for experience reuse (\Cref{def:projection}).
% \end{enumerate}

% \begin{definition}[Forward Model]\label{def:forward-model}

Given an initial configuration $x_k \in \mathcal{X}$ and a set of feasible control parameters $u \in \mathcal{U}_{MP}$, the \textsl{forward model} $f_{MP}$ for a motion primitive $MP$ returns the final configuration $x_{k+1} \in \mathcal{X}$, i.e.,

\begin{center}
$x_{k+1} = f_{MP}(x_k, u)$
\end{center}
\end{definition}

where $\mathcal{U}_{MP}$ is the set of feasible control parameters for the chosen motion primitive $MP$.

% \begin{definition}[Control Mechanism]\label{def:control-mechanism}
Given a vector of control parameters $u \in \mathcal{U}_{MP}$ from a set of feasible parameters corresponding to a motion primitive $MP$, the \textsl{control mechanism} $\phi_{MP}$ generates a robot trajectory that executes $u$.
\begin{center}
$\phi_{MP} \longrightarrow \tau: [0, 1]$ s.t. $\tau(0) = q_s, \hspace{3pt} \tau(1) = q_f, \hspace{3pt} \dot{q}_s = \dot{q}_f = 0$
\end{center}
where $q_s, q_f \in \mathcal{R}$ correspond to the static start and end joint configurations of the robot.
\end{definition}


In this work, we leverage \textsl{poking} as our motion primitives for \gm{give a reason}. A discussion of its characteristics is provided next.